\chapter{Arhitectura sistemului}

\section{Vedere de ansamblu}

Arhitectura Trainee este distribuita si impartita clar pe straturi: aplicatie mobila, backend API, baza de date relationala si servicii externe. Structura monorepo pastreaza codul intr-un singur spatiu de lucru, dar fara sa amestece responsabilitatile dintre frontend si backend.

\begin{figure}[H]
\centering
\fbox{\parbox{0.9\textwidth}{
\textbf{Placeholder diagrama de context} \\
Nivel 1: client mobil (Expo React Native) \\
Nivel 2: API server (Express + TypeScript) \\
Nivel 3: PostgreSQL + PostGIS \\
Servicii externe: Stripe, RevenueCat, SMTP, AWS S3
}}
\caption{Diagrama de context a sistemului Trainee}
\end{figure}

\section{Arhitectura frontend}

Frontend-ul este impartit in module functionale:
\begin{itemize}[leftmargin=1.2cm]
    \item \texttt{app/} pentru ecrane si rute Expo Router;
    \item \texttt{features/} pentru logica pe domenii (auth, billing, gym, schedule, trainer, support);
    \item \texttt{src/api/} pentru comunicatia comuna cu backend-ul;
    \item \texttt{src/constants/} pentru configurari de runtime.
\end{itemize}

Aceasta organizare limiteaza cuplarea dintre componente si permite extindere graduala. Comunicatia cu backend-ul este standardizata prin RTK Query, ceea ce aduce cache, invalidare si retry intr-o forma controlabila.

\subsection{Fluxul de date in frontend}

Fluxul uzual este direct:
\begin{enumerate}[leftmargin=1.2cm]
    \item utilizatorul declanseaza o actiune in UI;
    \item ecranul apeleaza un hook RTK Query;
    \item request-ul HTTP pleaca spre backend, cu token daca sesiunea este autentificata;
    \item starea se actualizeaza in store, iar interfata se rerandeaza.
\end{enumerate}

In practica, acest tip de flux face comportamentul mai previzibil si simplifica tratarea erorilor de retea.

\section{Arhitectura backend}

Backend-ul foloseste Express si o separare clasica pe layere:
\begin{itemize}[leftmargin=1.2cm]
    \item \textbf{Routes}: definirea endpoint-urilor pe domenii;
    \item \textbf{Middleware}: validare, autentificare, autorizare, rate limiting;
    \item \textbf{Controllers}: orchestration pentru cazurile de utilizare;
    \item \textbf{Models}: mapare ORM pe tabelele PostgreSQL;
    \item \textbf{Services}: integrare externa si utilitare de domeniu.
\end{itemize}

\subsection{Pipeline HTTP}

Pe endpoint-urile principale, ordinea executiei este:
\begin{enumerate}[leftmargin=1.2cm]
    \item parsare request;
    \item aplicare profile de rate limit;
    \item validare schema pe \texttt{body/query/params};
    \item autentificare JWT, daca ruta este protejata;
    \item autorizare pe rol, unde este cazul;
    \item executie controller;
    \item raspuns standardizat.
\end{enumerate}

Acest pipeline separa clar controalele defensive de logica de business si reduce duplicarea de cod.

\section{Persistenta si modelul relational}

Persistenta este realizata in PostgreSQL. Entitatea centrala este \texttt{User}, extinsa cu profil de antrenor, relatii cu sali, sloturi de programare si evenimente operationale.

Zonele cu impact major sunt:
\begin{itemize}[leftmargin=1.2cm]
    \item date de identitate si roluri;
    \item profiluri de antrenori, specializari si recenzii;
    \item programari, coduri de check-in si statusuri de slot;
    \item issue management si evenimente billing pentru idempotenta.
\end{itemize}

PostGIS si pg\_trgm completeaza modelul pentru cautare geografica si textuala la cost operational rezonabil.

\section{Integrarea cu servicii externe}

\subsection{Billing}

Fluxurile de billing sunt gandite sa functioneze in moduri de runtime diferite. Endpoint-urile de webhook sunt separate si procesate idempotent, tocmai pentru a evita efecte secundare in scenarii de retry.

\subsection{Email}

Email-ul este folosit pentru verificare cont si resetare parola. Separarea in servicii dedicate izoleaza dependinta de furnizor si simplifica inlocuirea ulterioara.

\subsection{Storage media}

Resursele media sunt stocate in S3, iar backend-ul gestioneaza metadatele si asocierea cu entitatile de domeniu relevante.

\section{Arhitectura de securitate}

Modelul de securitate este stratificat:
\begin{itemize}[leftmargin=1.2cm]
    \item \textbf{Layer 1 - Input hardening}: validare stricta si control pentru campuri neasteptate;
    \item \textbf{Layer 2 - Traffic shaping}: rate limiting profilat;
    \item \textbf{Layer 3 - Identity and access}: JWT + autorizare pe rol;
    \item \textbf{Layer 4 - Secret management}: secrete obligatorii in configurare;
    \item \textbf{Layer 5 - Business anti-abuse}: reguli anti-inflatie pentru vizualizari profil.
\end{itemize}

\begin{figure}[H]
\centering
\fbox{\parbox{0.9\textwidth}{
\textbf{Placeholder diagrama securitate multi-layer} \\
Input Validation -> Rate Limiting -> Auth/Role -> Business Rules -> Persistence
}}
\caption{Model defensiv pe mai multe straturi}
\end{figure}

\section{Scalabilitate si extensibilitate}

Desi proiectul este in etapa academica, structura actuala sustine extensii reale:
\begin{itemize}[leftmargin=1.2cm]
    \item inlocuirea rate limiter in-memory cu un store distribuit;
    \item introducerea observabilitatii avansate (Prometheus + Grafana);
    \item separarea unor domenii in servicii dedicate, daca traficul o cere;
    \item adaugarea de metode noi de autentificare si plata.
\end{itemize}

In ansamblu, arhitectura aleasa echilibreaza viteza de livrare cu cerintele de siguranta si mentenabilitate.
